% !TeX root = thesis.tex

\chapter{Evaluation}
\label{ch:evaluation}

This chapter evaluates the two main outcomes of the project: the quality of
the automatically generated question--answer (QA) annotations and their use as
a benchmark for vision-language models (VLMs). The experiments address three
questions: whether the generated annotations are reliable, whether the
benchmark separates model families, and how performance changes with visual
modality, model size, and input representation.


\section{Experimental Setup}
\label{sec:evaluation_setup}

\subsection{Dataset and Quality Control}

The annotation pipeline processes aligned RGB, infrared, event, and
depth recordings. Its initial output contained 2,730 QA blocks. Because many
blocks contained several numbered questions, preprocessing split them into
7,310 atomic candidates. Nine candidates with irrecoverable question--answer
count mismatches were removed. Rule-based checks and an LLM-assisted semantic
verifier then assessed answerability, answer matching, caption support,
modality appropriateness, single-question structure, and hallucination
risk. This process retained 5,465 valid visual QA items that form the VLM benchmark
(Table~\ref{tab:dataset_summary}).

Each of the 5,465 visual benchmark questions has one source modality. The
benchmark does not contain questions that require two or more visual
modalities to be jointly supplied and reasoned over. A separately generated
cross-modal disambiguation collection is therefore outside the scope of the
main benchmark results reported in this chapter.

\begin{table}[ht]
    \centering
    \caption{Final visual benchmark composition.}
    \label{tab:dataset_summary}
    \begin{tabular}{lr}
        \toprule
        \textbf{Modality} & \textbf{Valid QA items} \\
        \midrule
        RGB      & 2,318 \\
        Infrared & 1,561 \\
        Event    &   822 \\
        Depth    &   764 \\
        \midrule
        Total    & 5,465 \\
        \bottomrule
    \end{tabular}
\end{table}


\subsection{Benchmark Protocols and Models}

The source-modality baseline answers every visual question using the modality
from which it was generated. Each model receives the same eight
deterministically selected frames: four from the day or with-light recording
and four from the aligned night or no-light recording. Sampling is disabled,
the batch size is one, and answers are limited to 128 new tokens.

A primary modality-replacement evaluation uses 5,331 questions for which all four visual
representations are available. Every question is answered independently from
RGB, infrared, event, and depth input, resulting in 21,324 answers per model.
This measures robustness to replacing the visual representation; it does not
evaluate joint multimodal fusion. A final native-video experiment supplies the
complete aligned videos through each model's own video adapter. Since decoding
and tokenization differ between model families, this experiment compares
complete model--adapter systems rather than identical visual evidence.

We evaluate Qwen3-VL, InternVL3.5, and Molmo2 at approximately 4B and 8B
parameters. Failed or empty outputs remain in the result set. All configurations
use the same questions, prompts, and failure policy.


\subsection{Metrics}

The primary metric is a blinded semantic judge that labels an answer as
correct, partially correct, incorrect, or unjudgeable. Strict accuracy counts
only correct answers, while soft accuracy assigns half credit to partially
correct answers:

\[
\mathrm{Acc}_{\mathrm{strict}} =
\frac{N_{\mathrm{correct}}}
{N_{\mathrm{correct}}+N_{\mathrm{partial}}+N_{\mathrm{incorrect}}},
\qquad
\mathrm{Acc}_{\mathrm{soft}} =
\frac{N_{\mathrm{correct}}+0.5N_{\mathrm{partial}}}
{N_{\mathrm{correct}}+N_{\mathrm{partial}}+N_{\mathrm{incorrect}}}.
\]

Unjudgeable items are excluded from both denominators. A reproducible
task-aware score complements the judge: it uses Boolean accuracy for yes/no
questions, numeric accuracy for counting, ANLS for OCR, sequence-order scoring
for temporal tasks, set F1 for multi-concept answers, and token F1 otherwise.
To keep the main evaluation concise, the result tables report strict judge
accuracy as the primary metric. Soft judge accuracy, task-aware scores, token
F1, latency, and peak GPU memory were retained as supporting diagnostics but
are omitted from the main tables. They can be included in supplementary
material if required.


\section{Annotation Quality}
\label{sec:annotation_quality}

The raw output required substantial post-processing. Rule-based analysis found
that 1,263 of the original 2,730 blocks passed directly, 614 required
transformation, and 853 contained quality warnings. After splitting, the
semantic verifier passed 5,916 of 7,301 atomic candidates, marked 507 for
review, and rejected 878. Caption answerability and caption support were the
weakest dimensions, whereas answer consistency and modality appropriateness
both exceeded 98\%.

\begin{table}[ht]
    \centering
    \caption{Summary of annotation quality control.}
    \label{tab:annotation_quality}
    \begin{tabular}{llrr}
        \toprule
        \textbf{Stage} & \textbf{Outcome} & \textbf{Count} & \textbf{Rate} \\
        \midrule
        Rule-based blocks & Pass & 1,263 & 46.3\% \\
                          & Transform & 614 & 22.5\% \\
                          & Warning & 853 & 31.2\% \\
        \midrule
        Semantic candidates & Pass & 5,916 & 81.0\% \\
                            & Review & 507 & 6.9\% \\
                            & Reject & 878 & 12.0\% \\
        \midrule
        Final candidates & Retained & 5,465 & 74.8\% \\
        \bottomrule
    \end{tabular}
\end{table}

The most frequent structural issue was several questions appearing in one
generated block: 5,185 atomic items originated from such blocks. Other checks
found explicit frame references, unusually short answers, duplicates, and
unclear responses. These results show that splitting and semantic verification
are essential pipeline stages rather than optional cleanup. The final 5,465
retained visual items prevent malformed or unsupported annotations from
dominating the benchmark.


\section{Benchmark Results}
\label{sec:benchmark_results}

\subsection{Primary Modality-Replacement Evaluation}
\label{sec:primary_modality_replacement}

The primary model-capability analysis uses the same-question
modality-replacement benchmark. It contains 5,331 underlying questions for
which aligned RGB, infrared, event, and depth representations are available.
Each question is evaluated independently with all four input modalities,
producing 21,324 question--modality evaluations per model.

This paired and modality-balanced design controls for question content and
allows a more direct comparison of robustness across visual representations
than the source-modality baseline, whose QA counts differ by modality. The
21,324 evaluations are repeated measurements of 5,331 underlying questions,
not 21,324 fully independent QA samples. Each inference still contains only
one visual modality; the experiment evaluates modality replacement rather
than joint cross-modal reasoning.

Replacing the source modality can reduce accuracy because some evidence is not
preserved across RGB, infrared, event, and depth. For example, depth suppresses
color and text, while event data emphasizes motion rather than static
appearance. Nevertheless, Qwen3-VL remains strongest at both scales
(Table~\ref{tab:cross_modality_results}).

\begin{table*}[ht]
    \centering
    \caption{Strict accuracy on the primary same-question modality-replacement benchmark. Each score evaluates the same 5,331 underlying questions under four visual modalities. The modality breakdown uses eight fixed frames.}
    \label{tab:cross_modality_results}
    \begin{tabular}{llrrrrrr}
        \toprule
        \textbf{Family} & \textbf{Size} & \textbf{Video Overall} & \textbf{Frame Overall} & \textbf{RGB} & \textbf{Infrared} & \textbf{Event} & \textbf{Depth} \\
        \midrule
        InternVL3.5 & 4B & 0.4422 & 0.4362 & 0.4590 & 0.4671 & 0.4095 & 0.4091 \\
        Molmo2      & 4B & 0.4475 & 0.4282 & 0.5001 & 0.4828 & 0.4161 & 0.3137 \\
        Qwen3-VL    & 4B & 0.4792 & 0.5104 & 0.5825 & 0.5684 & 0.4908 & 0.3999 \\
        \midrule
        InternVL3.5 & 8B & 0.4420 & 0.4801 & 0.5136 & 0.5192 & 0.4476 & 0.4400 \\
        Molmo2      & 8B & 0.3986 & 0.4401 & 0.5070 & 0.4838 & 0.4196 & 0.3498 \\
        Qwen3-VL    & 8B & 0.4870 & \textbf{0.5404} & \textbf{0.5973} & \textbf{0.5739} & \textbf{0.5358} & \textbf{0.4547} \\
        \bottomrule
    \end{tabular}
\end{table*}

Under fixed-frame input, scaling from 4B to 8B improves all three model
families. The same pattern does not hold for native video: InternVL3.5 remains
approximately unchanged, Molmo2 decreases, and Qwen3-VL improves only
slightly. Fixed frames also outperform native video for all three 8B models.
These results indicate that modality robustness depends on the video adapter
and evidence representation rather than on parameter count alone. As illustrated
in Figure~\ref{fig:modality_heatmap}, performance degrades significantly across all models
when evaluating on Event and Depth modalities compared to RGB and Infrared.

\begin{figure}[!htbp]
    \centering
    \includegraphics[width=\textwidth]{images/02_modality_performance_heatmap.png}
    \caption{Heatmap of modality-specific performance across different models and sizes. It illustrates that all models exhibit performance drops on Event and Depth modalities compared to RGB and Infrared.}
    \label{fig:modality_heatmap}
\end{figure}

The difference between modalities is also task dependent. RGB and infrared are
generally more suitable for appearance, text, and object identity, while depth
better preserves geometry and event data can highlight motion boundaries.
Consequently, the benchmark should be interpreted using both the aggregate
ranking and modality-specific results.


\subsection{Source-Modality Baseline}
\label{sec:source_modality_baseline}

The source-modality baseline evaluates each of the 5,465 visual QA items using
the modality from which the question was generated. In contrast to the
modality-replacement experiment, this setting preserves the intended evidence
source for each question. Its purpose is to verify that the retained QA items
remain answerable under their intended modality and that the generated
benchmark can distinguish between model families.

Table~\ref{tab:matched_results} presents the controlled eight-frame results.
All models answered all 5,465 questions. Qwen3-VL achieves the best accuracy at
both scales: its 4B variant reaches 55.63\% strict accuracy and already
outperforms the 8B variants of InternVL3.5 and Molmo2. Qwen3-VL-8B obtains the
best overall result at 57.49\%.

\begin{table}[ht]
    \centering
    \caption{Strict accuracy on the source-modality baseline with eight fixed frames.}
    \label{tab:matched_results}
    \begin{tabular}{lrrr}
        \toprule
        \textbf{Family} & \textbf{4B} & \textbf{8B} & \textbf{Difference} \\
        \midrule
        InternVL3.5 & 0.4611 & 0.4942 & +0.0331 \\
        Molmo2      & 0.4516 & 0.4878 & +0.0362 \\
        Qwen3-VL    & \textbf{0.5563} & \textbf{0.5749} & +0.0186 \\
        \bottomrule
    \end{tabular}
\end{table}

Model family has a larger effect than parameter count alone. Scaling from 4B
to 8B improves strict accuracy by only 1.87--3.62 percentage points in this
setting. Paired McNemar tests show no significant difference between the 4B
InternVL3.5 and Molmo2 models ($p=0.250$), whereas Qwen3-VL significantly
outperforms both ($p<10^{-38}$).


\subsection{Source-Modality Input and Scale Sensitivity}

Native video does not consistently improve performance
(Table~\ref{tab:input_comparison}). It benefits InternVL3.5-4B and Molmo2-4B,
but Qwen3-VL performs better with the controlled frames. All three 8B models
are equal or weaker with their native-video adapters. This does not establish
that video is intrinsically inferior; rather, results depend on model-specific
frame decoding, tokenization, and context allocation. Figure~\ref{fig:scaling_effect}
illustrates these model scaling effects (4B to 8B) across both fixed-frame and
native-video input representations.

\begin{table}[ht]
    \centering
    \caption{Strict accuracy for fixed eight-frame and native-video input on
    the source-modality baseline. Each question is evaluated using its
    original evidence modality.}
    \label{tab:input_comparison}
    \begin{tabular}{llrrr}
        \toprule
        \textbf{Family} & \textbf{Size} & \textbf{Frames} &
        \textbf{Video} & \textbf{Difference} \\
        \midrule
        InternVL3.5 & 4B & 0.4611 & 0.4898 & +0.0287 \\
        Molmo2      & 4B & 0.4516 & 0.4966 & +0.0450 \\
        Qwen3-VL    & 4B & 0.5563 & 0.5360 & -0.0203 \\
        InternVL3.5 & 8B & 0.4942 & 0.4897 & -0.0045 \\
        Molmo2      & 8B & 0.4878 & 0.4727 & -0.0151 \\
        Qwen3-VL    & 8B & 0.5749 & 0.5420 & -0.0329 \\
        \bottomrule
    \end{tabular}
\end{table}

\begin{figure}[!htbp]
    \centering
    \includegraphics[width=0.8\textwidth]{images/03_model_scaling_effect.png}
    \caption{Model scaling effects (4B to 8B) under fixed-frame and native-video input representations. Native video does not consistently benefit from increased parameter count.}
    \label{fig:scaling_effect}
\end{figure}

A supplementary 4B experiment increased the frame count from eight to thirty.
InternVL improved from 46.11\% to 51.70\% strict accuracy, while Molmo2 and
Qwen3-VL improved by only 0.23 and 0.55 percentage points. More frames therefore
help only when the architecture can use the additional visual tokens
effectively. The InternVL run used a slightly different recorded model version
and should be treated as a sensitivity analysis rather than a strict ablation.
Figure~\ref{fig:efficiency_tradeoff} depicts the accuracy--efficiency trade-off
across the evaluated models, highlighting the balance between accuracy and computational resources.

\begin{figure}[!htbp]
    \centering
    \includegraphics[width=0.8\textwidth]{images/05_accuracy_efficiency_tradeoff.png}
    \caption{Accuracy--efficiency trade-off across evaluated models. Qwen3-VL-4B offers a particularly strong balance between strict accuracy and resource requirements.}
    \label{fig:efficiency_tradeoff}
\end{figure}


\section{Discussion and Limitations}
\label{sec:evaluation_discussion}

The experiments lead to three main findings. First, automatic annotation
requires both structural and semantic quality control. Second, the resulting
benchmark clearly separates model families, with Qwen3-VL providing the best
overall strict accuracy. Third, neither a larger language
model nor more visual input guarantees higher accuracy; the visual adapter and
the selection of relevant evidence are equally important.

Per-item error records support these conclusions. Frequent failures include
incorrect counts, reversed temporal order, hallucinated objects, incomplete
spatial relations, and confusion between day and night recordings. Some
modality-replacement failures are unavoidable because the replacement
modality does not contain the information required by the question. These cases make the
benchmark useful for analyzing modality dependence, but they should not be
interpreted as ordinary matched-input failures.

Several limitations remain. Most reference annotations are generated
automatically, and an LLM participates in generation, verification, and answer
judging, which may introduce evaluator bias. A stratified human study would
provide stronger validation. The cross-modality experiment evaluates one
modality at a time rather than multimodal fusion, and native-video adapters do
not expose identical evidence across model families. Consequently, the main
benchmark cannot support claims about successful joint cross-modal reasoning;
such claims would require evaluating the separate cross-modal disambiguation
QA collection under a dedicated multi-input protocol. Finally, the modalities
and task categories are imbalanced. Although accuracy remains comparable under
the fixed manifest, efficiency metrics would require separate interpretation
because the two model-size groups were run on different hardware.


\section{Summary}
\label{sec:evaluation_summary}

The pipeline converts 2,730 generated QA blocks into 7,301 cleaned atomic
candidates and retains 5,465 valid visual QA items. On the controlled benchmark, Qwen3-VL-4B and Qwen3-VL-8B reach 55.63\% and
57.49\% strict accuracy, respectively. The modality-replacement and video
experiments show that performance depends strongly on how visual evidence is
represented; they do not evaluate joint multimodal reasoning. Increasing model
size, frame count, or video duration alone is not sufficient for consistent
improvement.
